%HOMEWORK template for M 325K
\documentclass{article}
\headheight=7pt         \topmargin=14pt
\textheight=574pt       \textwidth=445pt
\oddsidemargin=18pt     \evensidemargin=18pt

%Begin package import

\usepackage{amsmath, amssymb, amsthm}
\usepackage{mathtools}
\usepackage{pdfpages}
\usepackage{tikz-cd}

%End package import
%Begin custom commands

\newcommand{\C}{\mathbb{C}}
\newcommand{\R}{\mathbb{R}}
\newcommand{\Q}{\mathbb{Q}}
\newcommand{\Z}{\mathbb{Z}}
\newcommand{\N}{\mathbb{N}}
\newcommand{\abs}[1]{\lvert #1\rvert}
\newcommand{\RM}[1]{\mathrm{#1}}
\newcommand{\BF}[1]{\textbf{#1}}
\newcommand{\e}{\epsilon}
\newcommand{\ceil}[1]{\lceil{#1}\rceil}
\newcommand{\floor}[1]{\lfloor{#1}\rfloor}
\newcommand{\rarr}[0]{\rightarrow}
\newcommand{\IM}[1]{\text{Im}(#1)}
\newcommand{\Hom}[0]{\text{Hom}}
\newcommand{\Pow}[1]{\text{Pow}(#1)}
\newcommand{\angles}[1]{{\langle #1 \rangle}}

%End custom commands
%Begin custom theorems

\newtheorem{innercustomgeneric}{\customgenericname}
\providecommand{\customgenericname}{}
\newcommand{\newcustomtheorem}[2]{%
\newenvironment{#1}[1]
{%
\renewcommand\customgenericname{#2}%
\renewcommand\theinnercustomgeneric{##1}%
\innercustomgeneric%
}
{\endinnercustomgeneric}
}

\newcustomtheorem{THRM}{Theorem}
\newcustomtheorem{LEM}{Lemma}
\newcustomtheorem{EX}{Exercise}
\newcustomtheorem{COR}{Corollary}
\newcustomtheorem{PROB}{Problem}
\newcustomtheorem{claim}{Claim}

\newenvironment{solution}
{\renewcommand\qedsymbol{$\blacksquare$}\begin{proof}[Solution]}
{\end{proof}}

\newenvironment{definition}
{\renewcommand\qedsymbol{$\blacksquare$}\begin{proof}[Definition]}
{\end{proof}}

%End custom theorems
\begin{document}

\title{Representation Theory 1}
\author{Arham Lodha}%put your name here
\date{January 28, 2024}%enter the due date here
\maketitle

\begin{definition}
    A linear representation of a group G means a vector space V and a homomorphism $G \to GL(V)$ (invertible linear maps from V to itself, a group under composition). 
    We often say G acts linearly on V.  We usually leave the name of $G \to GL(V)$ out of the notation.
\end{definition}

\bigskip

\begin{PROB}{1}\label{Problem 1.}
    Why do we call this a "representation" of G? We say the representation is "faithful" if $G \to GL(V)$ is injective. Why do we call this a "faithful representation"?
\end{PROB}
\begin{solution}
    This is called a representation of G because inside $GL(V)$, more specifically the image of the homomorphism, all of the generators and relations are captured. We say that the representation is faithful if inside $GL(V)$, more specifically the image of the homomorphism, if G is naturally inside $GL(V)$ .  
\end{solution}

\bigskip

\begin{PROB}{2a}\label{Problem 2a.}
    Lets say G acts on V and W. What should mean by $Hom_G(V,W)$ -- a map of G reps?
\end{PROB}
\begin{proof}
    $\forall f \in Hom_G(V, W)$, $f \in Hom_F(V, W)$ where $\forall g \in G$, $\forall v \in V$, $f(g \cdot v) = g \cdot f(v)$.    
\end{proof}

\begin{PROB}{2b}\label{Problem 2b.}
    Say we have two n-dim reps of $G, f_i: G \to GL_n(F)$, $i = 1, 2$. When are they isomorphic? Concretely, what does it mean about the matrices $f_1(g), f_2(g)$ over all g?
\end{PROB}
\begin{solution}
    When $f_1 \cong f_2$,  $\exists \alpha \in Hom_G(F^n, F^n)$, be the isomorphism of G-representations. Thus $\forall g \in G$, $\alpha \circ f_1(g) = f_2(g) \circ \alpha \implies \alpha \circ f_1(g) \circ \alpha^{-1} = f_2(g)$. Thus the matrices, $f_1(g)$ and $f_2(g)$ are conjugate.   
\end{solution}

\bigskip

\begin{PROB}{3a}\label{Problem 3.}
    Explain why there is a natural homomorphism $det: GL(V) \to F^*$, and a natural additive homomorphism Trace: $Hom_F(F,F) \to F$.  
\end{PROB}
\begin{solution}
    There is a natural homomorphism $\det: GL(V) \to F^*$ which factors through the isomorphism $\det: GL_n(F) \to F^*$ after the choice of basis. Similarly there is a natural additive homomorphism $tr : Hom_F(F^n, F^n) \to F$ which factors through the $tr: GL_n(F) \to F$ after the choice of basis. Trace and determinant is independent of choice of basis. 
\end{solution}

\begin{claim}{3a.i}\label{Claim 3a.i.}
    For $A, B \in GL_n(F)$, $tr(AB) = tr(BA)$
\end{claim}
\begin{proof}
    $tr(AB) = \sum_{i = 1}^{n} \sum_{j = 1}^{n} A_{i, j} B_{j, i}$ and $tr(BA) = \sum_{k = 1}^{n} \sum_{l = 1}^{n} B_{k, l} A_{l, k} = \sum_{l = 1}^{n} \sum_{k = 1}^{n} A_{l, k} B_{k, l} = \sum_{i = 1}^{n} \sum_{j = 1}^{n} A_{i, j} B_{j, i} = Tr(AB)$. Thus $tr(AB) = tr(BA)$
\end{proof}

\begin{claim}{3a.ii}\label{Claim 3a.ii.}
    For $A, B \in GL_n(F)$, $tr(BAB^{-1}) = tr(A)$
\end{claim}
\begin{proof}
    By claim 3a.i, $tr(BAB^{-1}) = tr(BB^{-1}A) = tr(IA) = tr(A)$. Thus $tr$ is independent of choice of basis.   
\end{proof}

\begin{claim}{3b}\label{Claim 3b.}
    If $f: G \to GL(V)$ is a rep, its "character", means the composition $\chi_f: tr \circ f: G \to F$. Show that if two reps are iso, they have identical characters (the character functions are equal).
\end{claim}
\begin{proof}
    Suppose $a: G \to GL(V)$ and $b: G \to GL(W)$ are representations of G which are isomorphic. Thus $\exists \alpha \in Hom_G(V, W)$ which is a isomorphism. By 2, we know $\forall g \in G$, $\alpha \circ a(g) = b(g) \circ \alpha \implies \alpha \circ a(g) \circ \alpha^{-1} = b(g)$, in other words isomorphic representations are conjugate, equivalently related by change of basis. $\forall g \in G$, $\chi_b(g) = tr(b(g)) = tr(\alpha \circ a(g) \circ \alpha^{-1})= tr(\alpha \circ \alpha^{-1} \circ a(g))$, by 3a.ii. Thus $\chi_b(g) =tr(b(g))= tr(\alpha \circ \alpha^{-1} \circ a(g)) = tr(a(g)) = \chi_a(g)$. Thus $\chi_b = \chi_a$.
\end{proof}

\begin{claim}{3c}\label{Claim 3c.}
    The amazing theorem of Frobenius, which we will prove, is that for finite dim complex representations, the converse is true: If two reps have the same characters, they are isomorphic. Why is this amazing (think in terms of preservation of information)?
\end{claim}
\begin{proof}
    When we take the character we lose information, it is amazing how regardless of the fact that we lost information we can tell that two reps are isomorphic.
\end{proof}


\bigskip

\begin{claim}{4}\label{Claim 4.}
    As a reality check, lets say G acts on V and W, and the characters are the same, and char(F) =0. Explain why V and W at least have the same dimensions.
\end{claim}
\begin{proof}
    Suppose $\alpha: G \to GL(V)$ and $\beta: G \to GL(W)$ are representations of G such that $\chi_\alpha : Tr \circ \alpha$ and $\chi_\beta : Tr \circ \beta$ and $\chi_\alpha = \chi_\beta$. Then $\alpha(e) = I_v \in GL(V)$ and $\beta(e) = I_w \in GL(W)$. The trace of identity is the dimension of the vector space, thus $tr(I_v) = \dim V $ and $tr(I_w) = \dim W$. Thus 
    
    \[
        \dim V = tr(I_v) = \chi_\alpha(e) = \chi_\beta(e) = tr(I_w) = \dim W.
    \]
\end{proof}

\bigskip

\begin{PROB}{5}\label{Problem 5.}
    Explain how we can "add" reps, ie if G acts on V and on W it naturally acts on the direct sum $V \times W$  (its the product set, but usually called the direct sum when we write + for the addition). Now $P := V \times W$ comes with two natural G invariant subspaces (iso to) $V, W \subset P$.  Now suppose G acts on (the F-vector space) P and we have two G invariant subspaces V, W. When is the natural addition map $V \oplus W \to P$ (sending $(v,w) \to v + w$, addition in P) is iso of G-reps? Find necessary and sufficient conditions. Given P, when we can find such V and W neither 0, we say P is "reducible". If its not, its irreducible.
\end{PROB}
\begin{proof}
    Let $\alpha: G \to GL(P)$. Let $V, W \leq P$ be G-invariant subspaces.   
    
    $\implies$: Suppose $V \cap W = 0$ and $\dim V + \dim W = \dim P$. 
    
    $\forall (v, w) \in \ker +$, $v + w = 0 \implies w = -v \implies v \in W$. Since $V \cap W = 0$, then $v = 0$ and $w = 0$. Thus $\ker + = (0,0)$ and $+$ is injective. 
    
    Let $v_1, \dots, v_{\dim V}$ be the basis of $V$ and $w_1, \dots, w_{\dim W}$ be the basis of $W$. Since $\ker + = 0$, $v_1, \dots, v_{\dim V}, w_1, \dots, w_{\dim W}$ are linearly independent. There are $\dim P$ vectors in $v_1, \dots, v_{\dim V}, w_1, \dots, w_{\dim W}$ which are linearly independent thus $v_1, \dots, v_{\dim V}, w_1, \dots, w_{\dim W}$ is a basis of P and thus $v_1, \dots, v_{\dim V}, w_1, \dots, w_{\dim W}$ span the space and $+$ is surjective and thus a isomorphism of vectorspaces.      
    
    Then $\forall v \in V$ and $\forall w \in W$, $+(g \cdot (v,w)) = +(\alpha(g)v, \alpha(g)w) = \alpha(g)v + \alpha(g)w = \alpha(g)(v + w) = \alpha(g)[+(v, w)]$. Thus $V \oplus W \to P$ is also a isomorphism of representations. 
    
    $\impliedby$: Suppose $V \oplus W \cong P$ as a g-rep. Then it is a isomorphism of vectorspaces. Then $\ker + = 0$, then $\forall v \in V$ and $\forall w \in W$, $v + w = 0 \implies v = 0$ and $w = 0$. $\forall v \in V \cap W$, $-v \in V \cap W$, $-v \in V \cap W$, $v - v = 0$. Thus $V \cap W \subset \ker +$. Thus $V \cap W = 0$. Then $\forall p \in P$, $p = v + w$, for some $v \in V$ and $w \in W$. Since $\ker + = 0$, V is linearly independent from W. Thus $\dim V + \dim W = \dim P$.
\end{proof}

\bigskip

\begin{PROB}{6}\label{Problem 6.}
    Prove any finite dim rep is a direct sum of irreducibles. Being reducible requires we have two G-invariant subspaces, which are complementary. The first theorem in the subject is that when F = C, we only need one (non-trivial) invariant subspace.
\end{PROB} 
\begin{proof}
    (Base Case): Let $V$, be a representation of $G$. Suppose $\dim V = 1$. Thus $V$  has no proper nontrivial subspaces, thus no proper G-invariant subspaces. Thus $V$ is a direct product of irreducibles, because it is irreducible itself. 
    
    (Induction Hypothesis): Let any representation of $G$ $V$ with $\dim V < k$ be a direct sum of irreducibles.
    
    (Inductive Step): Let $P$ be a representation of G with $\dim P = k$. If $P$ is irreducible, then we are done becuase $P$ itself is a direct product of irreducible, namely itself. If $P$ is reducible, by definition, there exists complementary G-invariant subspaces $V$ and $W$ such that there exists a g-rep isomorphism $V \oplus W \to P$. Thus $P \cong V \oplus W$. By definition of reducible, $V, W$ are nontrivial so $\dim V \geq 1$ and $\dim W \geq 1$. Thus $\dim V = \dim P - \dim W \leq k - 1$ and similarly $\dim W \leq k - 1$. Restrict the representation of $G$ to $V$ and $W$ and apply the inductive hypothesis. By the inductive hypothesis, $V \cong V_1 \oplus \dots \oplus V_m$ and $W \cong W_1 \oplus \dots \oplus W_n$ where $V_i$ for $i = 1, \dots, m$  and $W_j$ for $j = 1, \dots, n$ are irreducibles. Thus $P \cong V_1 \oplus \dots \oplus V_m \oplus W_1 \oplus \dots \oplus W_n$. Thus P is the direct product of irreducibles.                   
\end{proof}

\bigskip

\begin{PROB}{7}\label{Problem 7.}
    Suppose G acts on V, $\angles{ , } > 0$ (meaning its a positive def hermitian symmetric form) and $\angles{ , }$ is G-invariant, ie each g in G preserves $\angles{ , }$. Show that if W is G-invariant, then it has an invariant complement, and so V is reducible. 
\end{PROB}
\begin{proof}
    Let $W < V$, be a G-invariant subspace. $W^\perp := \left\{ v \in V : \angles{v, w} = 0, ~ \forall w \in W \right\}$. $\forall w \in W$ and $\forall g \in G$, since $W$ is G-invariant, $gw = w'$ for some $w' \in W$, so $w = g^{-1}w'$. Thus $\forall v \in W^\perp$,  $\angles{gv, w} = \angles{gv, g (g^{-1}w)}$. Since $\angles{ , }$ is G-invariant, $\angles{gv, g (g^{-1}w)} = \angles{v, g^{-1}w}$. Since $g^{-1}w \in W$, $\angles{gv, w} = angles{v, g^{-1}w} = 0$. So $gv \in W^\perp$. So $W^\perp$ is G-invariant. $\forall v \in W \cap W^\perp$, $v \in W$ and $v \in W^\perp$. Since $v \in W^\perp$, $\angles{v, v} = 0$. Since $\angles{ , }$ is positive definite, $v = 0$. Thus $W \cap W^\perp = 0$. Let $k = \dim W$ and let $w_1, \dots, w_k$ be a basis of W. $W^\perp = \ker\left(\left[v \to \begin{bmatrix}
        \angles{v, w_1} \\
        \vdots \\
        \angles{v, w_k}
    \end{bmatrix}\right]\right)$. Since $\angles{w_i, w_i} \neq 0$, so $w_i \to e_i$, so $\dim \left[v \to \begin{bmatrix}
        \angles{v, w_1} \\
        \vdots \\
        \angles{v, w_k}
    \end{bmatrix}\right]_*(V) = k$ (image of map). So by rank nullity, $\ker W^\perp = \dim V - k$. Thus $\dim W + \dim W^\perp = k + \dim V - k = \dim V$. Thus $W \oplus W^\perp \to V$ is a isomorphism. Thus $W^\perp$ is a G-invariant complement to V and  $V$ is reducible.                   
\end{proof}

\bigskip

\begin{PROB}{8a}\label{Problem 8a.}
    Define an action of $GL(V)$ on the set of pos def Herm symmetric forms  (I think there is only one way to do it).  Now let $\angles{ , } > 0$ be any positive form.   
\end{PROB}
\begin{solution}
    Let $P$ be the set of pos def Herm symmetric forms.
    $\forall g \in GL(V)$ and $\forall \angles{ , } \in P$, $\forall v, w \in V$, $\angles{v, w}_g := \angles{g^{-1}v , g^{-1}w}$. $g \cdot \angles{,} = \angles{v, w}_g$. $\forall v,w \in V$, $\forall g, h \in GL(V)$, and $\forall \angles{ , } \in P$, $(gh \cdot \angles{,})[v,w] = \angles{h^{-1}g^{-1}v, h^{-1}g^{-1}v} = \angles{,}_h[g^{-1}v, g^{-1}w] = (g \cdot \angles{,}_h)[v,w]= (g \cdot (h \cdot \angles{,}))[v, w]$.      
\end{solution}
\begin{claim}{8b}\label{Claim 8b.}
    Show that the sum of positive def Herm symmetric forms is also pos def.
\end{claim}
\begin{proof}
    Let $\angles{, }$ and $[ , ]$ be positive definite form hermitian symmetric forms. Let $p: V \times V \to F$ be a bilinear form where $\forall (v, w) \in V \times V$, $p(v, w) = \angles{v, w} + [v, w]$. $\forall v \in V \ni v \neq 0$, since $\angles{,}$ and $[,]$ are positive definite, $\angles{v,v} > 0$ and $[v,v] > 0$. So $p(v,v) = \angles{v, v} + [v, v] > 0$ since sum of 2 positive numbers is positive. Thus $p$ is positive definite.                
\end{proof}
\begin{PROB}{8c}\label{Problem 8c.}
    Let G be finite acting on V. Let [ , ] be the average over the orbit $G\cdot\angles{ , }$. Show this is G-invariant and positive.
\end{PROB}
\begin{proof}
    Let G be finite acting on V. Let $\angles{ , }$ be a positive hermitian symmetric bilinear form. $[,] := \frac{1}{|G|} \sum_{g \in G} g \cdot \angles{,}$. By 8b, since the sum of positive definite hermitian forms is positive, and $|G| > 0$, $[,] > 0$. $\forall h \in G$, $h \cdot [ , ]  = \frac{1}{|G|} h \cdot \sum_{g \in G} g \cdot \angles{,} = \frac{1}{|G|} \sum_{g \in G} hg \cdot \angles{,}$. $\forall g \in G$, $h h^{-1} g = g$, and since $h^{-1} \in G$, $h^{-1}g \in G$. Thus left multiplication of every element of G with $g$ simply reorders the group G. Thus, $h\cdot[,] = \frac{1}{|G|} \sum_{g \in G} g \cdot \angles{,} = \cdot[,]$. Thus $[,]$ is G-invariant.               
\end{proof}

\bigskip

\begin{PROB}{9}\label{Problem 9.}
    Show that any finite subgroup of $GL_n(C)$ is conjugate to a subgroup of $U(n)$ (the matrices that preserve the standard hermitian dot product).  Explain why every finite subgroup of $GL_n(R)$ is conjugate to a subgroup of O(n) (orthogonal matrices, matrices that preserve the standard dot product).
\end{PROB}
\begin{claim}{9a}\label{Claim 9a.}
    Let $P$ be the set of nondegenerate Herm symmetric forms. If $G := GL_n(\C)$ acts on $P$ using the action defined in 8, the action is transitive.

    $\forall v, w \in V$, $\angles{v, w}_g := \angles{g^{-1}v , g^{-1}w}$. $g \cdot \angles{,} = \angles{v, w}_g$.
\end{claim}
\begin{proof}
    Define $( , )$ as the standard hermitian dot product. Let $e_1, \dots. e_n$ be the standard basis.  
    
    $\forall \angles{ , } \in P$, by bilinear forms 2, $\exists$ a orthonormal basis $b_1, \dots, b_n \ni \angles{b_i, b_j} = \delta(i, j)$. Let $g \in G$, be the linear transformation which transforms the standard basis to the orthonormal basis. Thus $(g \cdot \angles{})[e_i, e_j] = \angles{e_i, e_j}_g = \angles{g^{-1}e_i, g^{-1}e_j} = \angles{b_i, b_j} = \delta(i, j) = (e_i, e_j)$. Thus $g \cdot \angles{} = (,)$. Thus all positive definite bilinear forms are in the same orbit as the standard hermitian dot product and thus are all in the same orbit and the action is transitive.   
\end{proof}
\begin{claim}{9b}\label{Claim 9b.}
    $U(n) = G^{(,)}$ (Stabilizer) 
\end{claim}
\begin{proof}
    $\forall g \in G^{(,)}$, $g \cdot (,) = (,)$. Let $h = g^{-1}$  Thus     

    \[
        [h^T\bar{h}]_{i,j} = e_i {h}^T \bar{h} \bar{e_j}  = (he_i, he_j) = (e_i, e_j)_g = (e_i, e_j) = \delta(i, j).
    \]
    Thus $[h^T\bar{h}]_{i,j} = \delta(i, j) \implies h^T\bar{h} = I \implies h \in U(n)$ and $g \in U(n)$ .

    $\forall g \in U(n)$, let $h = g^{-1}$ and $h \in U(n)$ . Thus $\forall v, w \in V$ 

    \[
        (g \cdot (,))[v, w] = (v, w)_g = (hv, hw) = v^T h^T \bar{h} \bar{w} = v^T \bar{w} = (v, w).
    \]

    Thus $g \cdot (,) = (,)$. 
    
    Thus $U(n) = G^{(,)}$ 
\end{proof}
\begin{claim}{9c}\label{Claim 9c.}
    Show that any finite subgroup of $H \leq GL_n(\C) = G$ is conjugate to a subgroup of $U(n)$ (the matrices that preserve the standard hermitian dot product).
\end{claim}
\begin{proof}
    Create the bilinear form $[ , ] = \frac{1}{|H|} \sum_{g \in H} g \cdot ( , )$. By 8, $[ , ]$ is $H$ invariant. So $H \leq G^{[ , ]}$. Since the action is transitive, $\exists T \in G = GL_n(\C) \ni T \cdot (, ) =[ , ]$. Thus $G^{[ , ]} = G^{T \cdot (, )} = T G^{(,)} T^{-1} = T U(n) T^{-1}$ (conjugation formula of stabilizers from last semester). Thus $H \leq T U(n) T^{-1}$.       
\end{proof}

\begin{claim}{9d}\label{Claim 9d.}
    Explain why every finite subgroup of $GL_n(R)$ is conjugate to a subgroup of O(n) (orthogonal matrices, matrices that preserve the standard dot product).
\end{claim}
\begin{proof}
    Using similar logic as above, we can prove this claim.

    Let $P$ be the set of positive definite symmetric bilinear forms. Let $G := GL_n(\R)$. 

    Define $( , )$ as the standard dot product. Let $e_1, \dots. e_n$ be the standard basis.

    $\forall \angles{ , } \in P$, by bilinear forms 2, $\exists$ a orthonormal basis $b_1, \dots, b_n \ni \angles{b_i, b_j} = \delta(i, j)$. Let $g \in G$, be the linear transformation which transforms the standard basis to the orthonormal basis. Thus $(g \cdot \angles{})[e_i, e_j] = \angles{e_i, e_j}_g = \angles{g^{-1}e_i, g^{-1}e_j} = \angles{b_i, b_j} = \delta(i, j) = (e_i, e_j)$. Thus $g \cdot \angles{} = (,)$. Thus all positive definite symmetric bilinear forms are in the same orbit as the standard dot product and thus are all in the same orbit and the action is transitive.


    $\forall g \in G^{(,)}$, $g \cdot (,) = (,)$. Let $h = g^{-1}$  Thus     

    \[
        [h^Th]_{i,j} = e_i {h}^T h e_j  = (he_i, he_j) = (e_i, e_j)_g = (e_i, e_j) = \delta(i, j).
    \]
    Thus $[h^Th]_{i,j} = \delta(i, j) \implies h^Th = I \implies h^{-1} = h^T \implies h \in O(n)$ and $g \in O(n)$ .

    $\forall g \in O(n)$, let $h = g^{-1}$ and $h \in O(n)$ . Thus $\forall v, w \in V$ 

    \[
        (g \cdot (,))[v, w] = (v, w)_g = (hv, hw) = v^T h^T hw = v^T w = (v, w).
    \]

    Thus $g \cdot (,) = (,)$. 
    
    Thus $O(n) = G^{(,)}$

    $\forall H \leq G$, 
    
    Create the bilinear form $[ , ] = \frac{1}{|H|} \sum_{g \in H} g \cdot ( , )$. By 8, $[ , ]$ is $H$ invariant. So $H \leq G^{[ , ]}$. Since the action is transitive, $\exists T \in G = GL_n(\R) \ni T \cdot (, ) =[ , ]$. Thus $G^{[ , ]} = G^{T \cdot (, )} = T G^{(,)} T^{-1} = T O(n) T^{-1}$ (conjugation formula of stabilizers from last semester). Thus $H \leq T O(n) T^{-1}$.
\end{proof}


\end{document}